\documentclass[11pt,a4paper]{article}
\usepackage[T1]{fontenc}
\usepackage[utf8]{inputenc}
\usepackage{lmodern}
\usepackage{microtype}
\usepackage[a4paper,margin=2.25cm]{geometry}
\usepackage{longtable,booktabs,array,calc}
\usepackage{fancyvrb}
\usepackage{xcolor}
\usepackage{hyperref}
\usepackage{xurl}
\usepackage{fancyhdr}
\usepackage{enumitem}
\usepackage{parskip}
\usepackage{titlesec}
\definecolor{ddtadark}{RGB}{30,38,48}
\definecolor{ddtablue}{RGB}{38,82,126}
\hypersetup{colorlinks=true,linkcolor=ddtablue,urlcolor=ddtablue,citecolor=ddtablue}
\pagestyle{fancy}
\fancyhf{}
\lhead{Documentation-Driven Threat Analysis}
\rhead{BA2 closure}
\cfoot{\thepage}
\setlength{\headheight}{14pt}
\setlist[itemize]{leftmargin=1.5em,itemsep=0.2em,topsep=0.25em}
\setlist[enumerate]{leftmargin=1.7em,itemsep=0.2em,topsep=0.25em}
\titleformat{\section}{\Large\bfseries\color{ddtadark}}{\thesection}{0.6em}{}
\titleformat{\subsection}{\large\bfseries\color{ddtadark}}{\thesubsection}{0.6em}{}
\providecommand{\tightlist}{\setlength{\itemsep}{0pt}\setlength{\parskip}{0pt}}
\setlength{\emergencystretch}{4em}
\hbadness=10000
\hfuzz=20pt
\sloppy
\begin{document}
\section{DDTA BA2 relation/action vocabulary - R1}\label{ddta-ba2-relationaction-vocabulary---r1}

\textbf{Status:} CLOSED BY BA2-T4 / ACCEPTED FOR CURRENT THESIS SCOPE

\textbf{Closure baseline reviewed:} \texttt{d16743a6417196ebf53840b1210a645e9dda4245}

\textbf{Identity dependency:} \texttt{BA1\_MINIMAL\_BAE\_IDENTITY\_ONTOLOGY\_R1.md} (\texttt{BAReferent\ +\ BAProposition})

\textbf{Historical derivation:} BA2-T1 structural lower bound -\textgreater{} BA2-T2 vocabulary pressure test -\textgreater{} BA2-T3 cross-corpus regression -\textgreater{} BA2-T4 closure review.

\subsection{1. Closure statement}\label{closure-statement}

BA2 closes the minimum methodology-neutral proposition/relation/action semantics required by the current DDTA thesis evidence.

The closure does \textbf{not} claim to define a universal systems-modeling language or an exhaustive verb ontology. It closes a governed Base Analysis semantic contract sufficient for the reviewed facial-access and order-fulfillment corpora and for the bounded method-neutral analysis-consumer pressure already performed.

The accepted proposition shape is:

\begin{verbatim}
BAProposition
|- semanticOperatorKey   1
|- participation         1..*
|    |- roleKey          1
|    `- term             1
|- polarity              1
`- scopedModifier        0..*    [condition / temporalScope only]
\end{verbatim}

\texttt{semanticOperatorKey}, \texttt{roleKey}, \texttt{polarity} and \texttt{scopedModifier} are semantic vocabulary/structure, not new BAE identity families.

A participation \texttt{term} is normally a \texttt{BAReferent}. A controlled typed local value is permitted only when the meaning does not require independent identity across propositions, projections or change. If it does, BA1 requires promotion to \texttt{BAReferent} identity.

\subsection{2. Accepted current-scope semantic operator registry}\label{accepted-current-scope-semantic-operator-registry}

The following thirteen operator concepts are accepted as the current-scope minimum:

\begin{verbatim}
transfer
produce
create
observe
transition
correlate
reference
dependOn
consumeService
realize
assignResponsibility
constrain
classify
\end{verbatim}

Each key has a stable methodology-neutral meaning. Human-readable wording, authoring verbs and synonyms are not semantic identity and remain downstream lexical concerns.

\subsubsection{2.1 Normative meanings}\label{normative-meanings}

\begin{longtable}[]{@{}
  >{\raggedright\arraybackslash}p{(\columnwidth - 2\tabcolsep) * \real{0.5000}}
  >{\raggedright\arraybackslash}p{(\columnwidth - 2\tabcolsep) * \real{0.5000}}@{}}
\toprule\noalign{}
\begin{minipage}[b]{\linewidth}\raggedright
Operator
\end{minipage} & \begin{minipage}[b]{\linewidth}\raggedright
Accepted methodology-neutral meaning
\end{minipage} \\
\midrule\noalign{}
\endhead
\bottomrule\noalign{}
\endlastfoot
\texttt{transfer} & Assert conveyance of content from a source to one or more destinations. \\
\texttt{produce} & Assert that an actor/capability makes one or more governed results or outputs available, optionally from explicit inputs. \\
\texttt{create} & Assert establishment of a new project-semantic item or event occurrence. \\
\texttt{observe} & Assert inspection/read/query of existing project state or meaning without asserting creation or state change by the observation itself. \\
\texttt{transition} & Assert governed state/lifecycle change of a subject. \\
\texttt{correlate} & Assert same-request/evaluation/operation/context identity matching required to bind items correctly. \\
\texttt{reference} & Assert an explicit reference from one project meaning to another without by itself implying correlation, dependency or ownership. \\
\texttt{dependOn} & Assert that one project meaning requires another meaning/result/capability/milestone as a prerequisite. \\
\texttt{consumeService} & Assert actual consumption/use of a capability or service without transferring its ownership or responsibility. \\
\texttt{realize} & Assert that a more concrete project meaning realizes/materializes an abstract project meaning. \\
\texttt{assignResponsibility} & Assert placement or denial of a governed responsibility/authority kind over a project-semantic scope. \\
\texttt{constrain} & Assert a reusable or independently queryable restriction on behavior, state, acceptance, applicability or other governed project meaning. \\
\texttt{classify} & Assert assignment of a reusable method-neutral semantic kind to a \texttt{BAReferent}. \\
\end{longtable}

\subsection{3. Accepted operator-scoped role contracts}\label{accepted-operator-scoped-role-contracts}

Role concepts may recur across operators, but role validity and cardinality are operator-scoped. A context-free global role contract is rejected.

\begin{longtable}[]{@{}
  >{\raggedright\arraybackslash}p{(\columnwidth - 6\tabcolsep) * \real{0.2500}}
  >{\raggedright\arraybackslash}p{(\columnwidth - 6\tabcolsep) * \real{0.2500}}
  >{\raggedright\arraybackslash}p{(\columnwidth - 6\tabcolsep) * \real{0.2500}}
  >{\raggedright\arraybackslash}p{(\columnwidth - 6\tabcolsep) * \real{0.2500}}@{}}
\toprule\noalign{}
\begin{minipage}[b]{\linewidth}\raggedright
Operator
\end{minipage} & \begin{minipage}[b]{\linewidth}\raggedright
Required roles
\end{minipage} & \begin{minipage}[b]{\linewidth}\raggedright
Optional roles
\end{minipage} & \begin{minipage}[b]{\linewidth}\raggedright
Accepted current-scope cardinality
\end{minipage} \\
\midrule\noalign{}
\endhead
\bottomrule\noalign{}
\endlastfoot
\texttt{transfer} & \texttt{source}, \texttt{destination}, \texttt{content} & none & source \texttt{1}; destination \texttt{1..*}; content \texttt{1..*} \\
\texttt{produce} & \texttt{actor}, \texttt{result} & \texttt{input} & actor \texttt{1}; result \texttt{1..*}; input \texttt{0..*} \\
\texttt{create} & \texttt{actor}, \texttt{created} & none & actor \texttt{1}; created \texttt{1..*} \\
\texttt{observe} & \texttt{actor}, \texttt{observed} & \texttt{result} & actor \texttt{1}; observed \texttt{1..*}; result \texttt{0..*} \\
\texttt{transition} & \texttt{subject}, \texttt{toState} & \texttt{actor}, \texttt{fromState} & subject \texttt{1}; toState \texttt{1}; actor \texttt{0..1}; fromState \texttt{0..1} \\
\texttt{correlate} & \texttt{correlatedItem}, \texttt{correlationContext} & none & correlatedItem \texttt{1..*}; correlationContext \texttt{1} \\
\texttt{reference} & \texttt{referencer}, \texttt{referenced} & none & referencer \texttt{1}; referenced \texttt{1..*} \\
\texttt{dependOn} & \texttt{dependent}, \texttt{prerequisite} & none & dependent \texttt{1}; prerequisite \texttt{1..*} \\
\texttt{consumeService} & \texttt{consumer}, \texttt{service} & \texttt{provider} & consumer \texttt{1..*}; service \texttt{1}; provider \texttt{0..1} \\
\texttt{realize} & \texttt{abstract}, \texttt{realization} & none & abstract \texttt{1}; realization \texttt{1..*} \\
{\footnotesize\texttt{assignResponsibility}} & {\footnotesize\texttt{responsibleParty}, \texttt{responsibilityScope}, \texttt{responsibilityKind}} & {\footnotesize none} & {\footnotesize responsibleParty \texttt{1..*}; responsibilityScope \texttt{1}; responsibilityKind \texttt{1}} \\
\texttt{constrain} & \texttt{constraintTarget}, \texttt{constraintValue} & none & constraintTarget \texttt{1}; constraintValue \texttt{1..*} \\
\texttt{classify} & \texttt{classifiedReferent}, \texttt{semanticKind} & none & classifiedReferent \texttt{1}; semanticKind \texttt{1..*} \\
\end{longtable}

These cardinalities are accepted for the current thesis scope. They are not claims that all conceivable future projects must use these maxima. A future governed corpus may reopen the smallest affected BA2 contract if it demonstrates that widening or splitting one role contract is required to preserve method-neutral project meaning.

\subsection{4. Responsibility and ownership boundary}\label{responsibility-and-ownership-boundary}

\texttt{ownOrManage} is not a separate operator.

Ownership, management and other authority modes are represented through \texttt{assignResponsibility} plus a controlled \texttt{responsibilityKind} and explicit polarity.

Conceptually:

\begin{verbatim}
operator: assignResponsibility
responsibleParty    -> project
responsibilityScope -> underlying transport infrastructure
responsibilityKind  -> ownership / management
polarity            -> negative
\end{verbatim}

This preserves the facial-access distinction that the project can consume transport while not owning or managing the underlying infrastructure:

\begin{verbatim}
consumeService != responsibility/ownership
\end{verbatim}

\subsection{5. Explicit polarity}\label{explicit-polarity}

Every proposition exposes normalized polarity:

\begin{verbatim}
polarity = affirmative | negative
\end{verbatim}

Polarity prevents proliferation of \texttt{notX} operator keys and keeps negation separate from operator identity.

A negative proposition does not eliminate the need for \texttt{constrain}. Normative prohibitions whose content is itself a governed reusable rule remain explicit constraint propositions.

\subsection{6. Scoped modifier lower bound}\label{scoped-modifier-lower-bound}

Only two embedded modifier kinds are accepted in the BA2 lower bound:

\begin{itemize}
\tightlist
\item
  \texttt{condition} - a proposition-local applicability/guard condition with no independent reuse or identity;
\item
  \texttt{temporalScope} - a proposition-local before/after/during/until scope with no independent reuse or identity.
\end{itemize}

A modifier may remain embedded only when all of the following hold:

\begin{enumerate}
\def\labelenumi{\arabic{enumi}.}
\tightlist
\item
  it changes only the interpretation/applicability of one proposition or participation binding;
\item
  it introduces no independent participant set;
\item
  it needs no independent assertion-level provenance/review/change identity;
\item
  it is not reused elsewhere as project meaning.
\end{enumerate}

If any condition fails, the meaning is represented through \texttt{BAReferent} and/or a separate \texttt{BAProposition}.

The following are therefore \textbf{not} default embedded modifiers when governed or analytically relevant:

\begin{itemize}
\tightlist
\item
  completion/failure semantics;
\item
  atomicity;
\item
  concurrency;
\item
  idempotency;
\item
  reusable applicability rules;
\item
  realization;
\item
  correlation;
\item
  responsibility/authority;
\item
  service consumption;
\item
  reusable security or other specialized constraints.
\end{itemize}

They are represented through explicit propositions, normally \texttt{constrain} or another accepted operator as appropriate.

\subsection{7. Classification contract}\label{classification-contract}

Method-neutral referent classification is represented through a normal proposition:

\begin{verbatim}
operator: classify
classifiedReferent -> <BAReferent>
semanticKind       -> <controlled method-neutral semantic key/value>
\end{verbatim}

BA2 closes the \textbf{registry contract}, not an exhaustive universal semantic-kind taxonomy.

A semantic-kind registry entry must have:

\begin{enumerate}
\def\labelenumi{\arabic{enumi}.}
\tightlist
\item
  a stable semantic key;
\item
  a normative method-neutral definition;
\item
  enough distinction for a consumer not to reconstruct the kind from raw prose;
\item
  no method-specific meaning such as STRIDE category, DFD element or tool class.
\end{enumerate}

Examples under current pressure include party/person, capability, information/resource, behavior/event, store and contract. These examples do not constitute a universal closed enumeration and do not become BA1 metaclasses.

Classification propositions may later receive provenance, status and change disposition in BA3 while the classified \texttt{BAReferent} identity remains stable.

\subsection{8. Semantic key versus lexical wording}\label{semantic-key-versus-lexical-wording}

BA2 closes this separation:

\begin{verbatim}
source lexical wording
        !=
stable BA semantic key
        !=
display label / synonym
\end{verbatim}

For example, \texttt{deliver}, \texttt{send}, \texttt{submit} and \texttt{emit} may map to \texttt{transfer} when they express conveyance. \texttt{reads} and \texttt{queries} may map to \texttt{observe}. \texttt{maps} need not create a \texttt{map} operator when the governed meaning is production of a normalized result from an input plus a constraint.

The document-to-semantic mapping belongs to BA3 derivation mechanics. Display labels and authoring synonyms belong to BA5. Neither changes the accepted BA2 semantic key merely because wording changes.

\subsection{9. General logical composition}\label{general-logical-composition}

Current evidence does not require a general-purpose logical-expression language inside BA2.

Compound governed clauses may derive to multiple \texttt{BAProposition} identities sharing the relevant \texttt{BAReferent} meanings. N-ary participation preserves the scope of each assertion; \texttt{condition}, \texttt{temporalScope}, polarity and explicit \texttt{constrain} propositions preserve the reviewed conditional, negative and normative rule pressure.

A future corpus must reopen BA2 if it demonstrates that a governed atomic assertion cannot be preserved without an explicit logical-composition construct and cannot be represented honestly by the accepted proposition shape plus separate reusable propositions.

\subsection{10. Current-scope completeness and extension rule}\label{current-scope-completeness-and-extension-rule}

The thirteen-key registry is \textbf{closed for current thesis scope}, but it is not a universal verb ontology.

A new operator key is justified only if concrete governed evidence demonstrates all of the following:

\begin{enumerate}
\def\labelenumi{\arabic{enumi}.}
\tightlist
\item
  the project meaning is methodology-neutral;
\item
  existing operators plus roles, polarity, constraints and local modifiers cannot preserve the distinction without semantic loss or raw-prose reconstruction;
\item
  the distinction is needed for reusable human/method consumption or change reasoning;
\item
  the distinction is not merely lexical, implementation-specific, document-layer-specific or analysis-method-specific.
\end{enumerate}

If those conditions are met, reopen only the affected BA2 registry/role contract. Do not broaden Base Analysis into a general systems-modeling language.

A role contract is reopened only when a concrete governed case cannot be represented faithfully under the accepted operator-scoped roles/cardinalities. A modifier kind is added only when it remains proposition-local under the modifier-admissibility rule and cannot be represented more honestly through an explicit proposition.

\subsection{11. Relationship to BA1 and later phases}\label{relationship-to-ba1-and-later-phases}

BA2 does not reopen BA1. \texttt{BAReferent} and \texttt{BAProposition} remain the only first-class semantic identity families.

BA2 does not define:

\begin{itemize}
\tightlist
\item
  source locators, baseline provenance or grounded/derived/diagnostic state - BA3;
\item
  accepted/stale/superseded assertion lifecycle and equivalence mechanics - BA3;
\item
  human or method-specific views/projections - BA4;
\item
  lexical labels, synonyms or optional authoring/extraction assistance - BA5;
\item
  complete holdout regression and final Base Analysis closure - BA6;
\item
  STRIDE/STRIDE-AI categories, AnalysisRecord, Finding/Common Finding or ThreatForge classes - downstream analysis/tooling work.
\end{itemize}

Governed project documentation remains project authority. Base Analysis remains the accepted/rebuildable analytical representation of that governed meaning.

\subsection{12. Camera example under the closed contract}\label{camera-example-under-the-closed-contract}

A facial-access branch can expose, conceptually:

\begin{verbatim}
transfer
  source      -> CameraSubsystem
  destination -> RecognitionProcessor
  content     -> RecognitionCapture

correlate
  correlatedItem     -> RecognitionCapture
  correlationContext -> RecognitionRequest

consumeService
  consumer -> project endpoint context
  service  -> LocalConnectivity

assignResponsibility [negative]
  responsibleParty    -> project
  responsibilityScope -> underlying transport
  responsibilityKind  -> ownership/management

realize
  abstract    -> LocalConnectivity
  realization -> WiredEthernet

constrain
  constraintTarget -> delivery behavior
  constraintValue  -> incomplete delivery != successful completion

constrain
  constraintTarget -> delivery behavior
  constraintValue  -> confidentiality / integrity / authorized provenance
\end{verbatim}

A later threat-method overlay may project these facts into its own constructs. The overlay does not redefine the shared Base Analysis semantics.

\subsection{13. Closure disposition}\label{closure-disposition}

\begin{verbatim}
BA2 proposition shape                         ACCEPTED
Stable methodology-neutral semantic key       ACCEPTED
Explicit role-bound n-ary participation        ACCEPTED
Current-scope 13-key operator registry         ACCEPTED
Operator-family facet in normative core        REJECTED
Operator-scoped role/cardinality contracts     ACCEPTED
Explicit polarity                              ACCEPTED
Embedded modifiers: condition/temporalScope    ACCEPTED
Generic/untyped modifier bag                   REJECTED
Governed reusable rule -> constrain             ACCEPTED
Classification-as-proposition                  ACCEPTED
Universal semantic-kind taxonomy               NOT REQUIRED
Semantic key != lexical wording                CLOSED
Third BA1 identity family                      NOT FORCED
BA1                                            REMAINS CLOSED
BA2                                            CLOSED BY BA2-T4
\end{verbatim}

\end{document}
